% Cote 119: offprints, reports, courses, CV and letters (typed and printed part): archive cote 119, site page 178 (PDF page 179).
% typescript. Draft compared: none.
% AI-transcribed by Claude Opus 5.5 (claude-opus-5-5, adaptive thinking, effort high), Message Batches API, batch r119_full (msgbatch_01DwZNm5Y7YsKBWyAN4sFs8w), 2026-11-13; not yet checked by eye.
\scanpage{119/178}
\mspage{- 17 -}

6.5. \uline{Théorème}.

Soit $(t,u) \to \delta(t,u)$ une application continue d'un rectangle
$a \leq t \leq b$ , $a' \leq u \leq b'$ dans l'ouvert $U$ et soit $\omega$ une forme différen-%
tielle fermée dans $U$. Alors il existe une application $f(t;u)$ continue dans
le rectangle telle que au voisinage de tout point $(t_0,u_0)$ du rectangle
il existe une primitive $F$ de $\omega$ vérifiant $F(\delta(t;u)) = f(t,u\unsure{)|=}$ .
(Une telle fonction $f$ sera appelée primitive de $\omega$ suivant l'application $\delta$).
Cette primitive est unique à l'addition d'une constante près.

\uline{Démonstration}. Il s'agit encore d'une démonstration facile mais technique
basée sur la compacité du rectangle. En utilisant celle-ci on obtient des
subdivisions de $[a,b]$ et $[a',b']$
\[ a = t_0 < t_1 < t_2 \ \dots < t_n = b \]
\[ a' = u_0 < u_1 < u_2 \ \dots < u_m = b' \]
telles que le produit $[t_i, t_{i+1}] \times [u_j, u_{j+1}]$ soit un rectangle appliqué
par $\delta$ dans un disque ouvert $U_{i,j}$ dans lequel existe une primitive $F_{i,j}$
de $\omega$ . Fixons $j$. Comme $U_{i,j} \cap U_{\unsure{ix1|i\times 1},j}$ est non vide et connexe on peut
ajouter à chaque $F_{i,j}$ ($i$ variable) une constante telle que $F_{i,j}$ et
$F_{i+1,j}$ coïncident dans $U_{i,j} \cap U_{i+1, j}$ (remarque suivant la définition 3.2)
On obtient alors pour $u \in [u_j, u_{j+1}]$ une fonction $f_j(t,u)$ telle que
pour tout $i$ on ait
\[ f_j(t,u= = F_{i,j} \ (\delta(t,u)) \ \text{lorsque} \ t \in [t_i, t_{i+1}] \]
Ainsi $f_j(t,u)$ est continue dans le rectangle
\[ a \leq t \leq b \qquad u_j \leq u \leq u_{j+1} \]
Soit $V_j = \bigcup_{i=0}^{n} U_{i,j}$
$V_j \cap V_{i,j+1}$ est non vide et connexe car pour tout $i$ . $U_{i,j} \cap U_{i,j+1}$ est
connexe et contient $\delta([t_i,t_{i+1}] \times \{u_{j+1}\})$ . $[\delta([a,b] \times \{u_{j+1}\}$ est connexe.

\margin{178}
